% ============================== Part C ==============================
\problem{C1}{One potential gives both mass and shear}

\begin{given}
A thin lens is described by the lensing potential $\psi(\bm\theta)$. Its second derivatives
($\psi_{ij} = \partial^2\psi/\partial\theta_i\partial\theta_j$) give the convergence (projected mass in units of the
critical density) and the two shear components:
\[
\kappa = \tfrac12(\psi_{11} + \psi_{22}),\qquad
\gamma_1 = \tfrac12(\psi_{11} - \psi_{22}),\qquad
\gamma_2 = \psi_{12}.
\]
\source{\cBSpsi.}
A singular isothermal sphere (SIS) has $\psi = \theta_E\,|\bm\theta|$, with Einstein radius $\theta_E$.
\source{\cBSsis.}
\end{given}

\need{nothing (calculus only).}

\begin{enumerate}
  \item Compute $\psi_{11}$, $\psi_{22}$ and $\psi_{12}$ for the SIS. Write $\theta_1 = \theta\cos\phi$,
    $\theta_2 = \theta\sin\phi$.
  \item Show that $\kappa = \theta_E/(2\theta)$ and that $\gamma \equiv \gamma_1 + i\gamma_2 = -\frac{\theta_E}{2\theta}\,e^{2i\phi}$.
  \item Why does knowing $\kappa$ everywhere fix $\gamma$ everywhere (up to what)?
\end{enumerate}

\solution{
(a) $\partial_1\psi = \theta_E\theta_1/\theta$. Then $\psi_{11} = \theta_E\theta_2^2/\theta^3$,
$\psi_{22} = \theta_E\theta_1^2/\theta^3$ and $\psi_{12} = -\theta_E\theta_1\theta_2/\theta^3$.
(b) $\kappa = \tfrac12\theta_E(\theta_1^2 + \theta_2^2)/\theta^3 = \theta_E/(2\theta)$.
$\gamma_1 = \theta_E(\theta_2^2 - \theta_1^2)/(2\theta^3) = -\tfrac{\theta_E}{2\theta}\cos2\phi$ and
$\gamma_2 = -\tfrac{\theta_E}{2\theta}\sin2\phi$, so $\gamma = -\tfrac{\theta_E}{2\theta}e^{2i\phi}$.
(c) $\kappa$ fixes $\nabla^2\psi$, which fixes $\psi$ up to a harmonic function (one with $\nabla^2 = 0$, such as a
constant gradient or a term $\propto \theta^2$ that is really a uniform sheet; see C3). The shear follows from $\psi$.
This link is the basis of every mass map (Part D).
}

% --------------------------------------------------------------------
\problem{C2}{Tangential and cross shear}

\begin{given}
About the lens centre, the tangential and cross components of the shear are
$\gamma_t = -\mathrm{Re}(\gamma\,e^{-2i\phi})$ and $\gamma_\times = -\mathrm{Im}(\gamma\,e^{-2i\phi})$.
\source{\cBStan.}
For any mass distribution, the tangential shear averaged around a circle of radius $\theta$ is
$\langle\gamma_t\rangle(\theta) = \bar\kappa(\le\theta) - \langle\kappa\rangle(\theta)$, where $\bar\kappa(\le\theta)$ is
the mean convergence inside the circle and $\langle\kappa\rangle(\theta)$ the mean on it. For an axially symmetric lens
the averages can be dropped.
\source{\cBSgtmean.}
\end{given}

\need{C1 (SIS $\kappa$ and $\gamma$).}

\begin{enumerate}
  \item Use C1(b) to show that the SIS has $\gamma_t = \theta_E/(2\theta)$ and $\gamma_\times = 0$.
  \item Check the result against $\bar\kappa(<\theta) - \kappa(\theta)$.
  \item A subhalo next to a main halo breaks the axial symmetry about the main-halo centre. What does it add to
    $\gamma_\times$ measured about the main centre, and why do our polar cells stack in $(\gamma_t, \gamma_\times)$?
\end{enumerate}

\solution{
(a) $\gamma e^{-2i\phi} = -\tfrac{\theta_E}{2\theta}$ is real and negative, so $\gamma_t = \theta_E/(2\theta)$ and
$\gamma_\times = 0$.
(b) $\bar\kappa(<\theta) = \frac{1}{\pi\theta^2}\int_0^\theta \frac{\theta_E}{2\theta'}\,2\pi\theta'\,d\theta' = \theta_E/\theta$,
so $\bar\kappa - \kappa = \theta_E/\theta - \theta_E/(2\theta) = \theta_E/(2\theta)$. It agrees.
(c) The subhalo's shear is tangential about \emph{its own} centre, so about the main centre it has both $\gamma_t$ and
$\gamma_\times$ parts. A non-zero $\gamma_\times$ about the main centre signals non-axisymmetric structure. In the
$(\gamma_t, \gamma_\times)$ frame the main halo's signal sits in one component, so departures from it stand out.
}

% --------------------------------------------------------------------
\problem{C3}{Reduced shear and the mass-sheet degeneracy}

\begin{given}
Galaxy shapes respond to the reduced shear $g = \gamma/(1-\kappa)$, not to $\gamma$.
\source{\cBSreduced.}
The transformation $\kappa \to \lambda\kappa + (1-\lambda)$, $\gamma \to \lambda\gamma$ (adding a uniform sheet of
mass and rescaling) leaves $g$ unchanged: the mass-sheet degeneracy.
\source{\cBSmsd.}
\end{given}

\need{C1.}

\begin{enumerate}
  \item Verify that $g$ is invariant under the transformation.
  \item For $\kappa = 0.05$ and $\gamma = 0.02$, compute $g$. How large is the difference from $\gamma$, in percent?
  \item What can shear data \emph{not} tell you about a $\kappa$ map?
\end{enumerate}

\solution{
(a) $g' = \lambda\gamma/(1 - \lambda\kappa - 1 + \lambda) = \lambda\gamma/[\lambda(1-\kappa)] = g$.
(b) $g = 0.02/0.95 = 0.02105$, 5\% above $\gamma$. That is negligible far out, but not near the cluster centre, where
$\kappa \sim 0.2$.
(c) The overall level (and scale) of $\kappa$: only its shape is fixed. A mass map's zero point is a convention.
}

% --------------------------------------------------------------------
\problem{C4}{Adding two halos}

\begin{given}
$\kappa$ is the projected mass, and $\psi$ obeys the linear Poisson equation $\nabla^2\psi = 2\kappa$. So the $\psi$,
$\kappa$ and $\gamma$ of two halos are the sums of their separate values.
\source{\cBSpoisson.}
\end{given}

\need{C1, C3.}

\begin{enumerate}
  \item Halo A gives $\kappa_A = 0.05$, $\gamma_A = (0.02, 0)$ at some point; halo B gives $\kappa_B = 0.03$,
    $\gamma_B = (0, 0.01)$. Compute the correct reduced shear of A+B.
  \item Compute the incorrect $g_A + g_B$. How big is the error?
  \item Why do our simulations add $\gamma$ and $\kappa$ over the halos first and form $g$ only at the end?
\end{enumerate}

\solution{
(a) $\gamma = (0.02, 0.01)$ and $\kappa = 0.08$: $g = (0.02174, 0.01087)$.
(b) \[g_A + g_B = (0.02/0.95,\ 0.01/0.97) = (0.02105,\ 0.01031).\] The error is $\approx 3$--5\%, as large as the effects
we test for.
(c) Because $\gamma$ and $\kappa$ superpose and $g$ does not.
}

% --------------------------------------------------------------------
\problem{C5}{Where the subhalo's signal is, and how strong it is}

\begin{given}
In our simulation (subhalo of $3\times10^{14}\,M_\odot$ at $z = 0.13$), the subhalo's extra reduced shear at distance $r$
from it is close to $|\Delta g| \approx A/r$ with $A \approx 0.009$\,Mpc. That is SIS-like; the true profile is NFW,
whose shear formulae are in \cWBnfw. The sources have density $n \approx 5\,{\rm Mpc}^{-2}$, and each measures each
shear component with noise $\sigma \approx 0.05$. For a signal of known shape, the best possible (matched-filter)
significance is $({\rm S/N})^2 = \sum_i |\Delta\vg_i|^2/\sigma^2$. \emph{Derived here} (numbers from our
\texttt{s3\_sub1} manifest).
\end{given}

\need{A2 (information adds).}

\begin{enumerate}
  \item Replace the sum by an integral over a uniform source density and show that
    $({\rm S/N})^2 = 2\pi n A^2\sigma^{-2}\ln(r_{\max}/r_{\min})$.
  \item Evaluate it for $r_{\min} = 0.1$\,Mpc and $r_{\max} = 3$\,Mpc.
  \item What fraction of $({\rm S/N})^2$ comes from within 1\,Mpc, and from within 0.5\,Mpc, of the subhalo?
  \item With each galaxy's full, anisotropic precision matrix in place of an isotropic $\sigma$, the same calculation
    gives ${\rm S/N} \approx 3.2$ for this galaxy set. Why is it larger?
\end{enumerate}

\solution{
(a) $\sum_i \to \int n\,2\pi r\,dr$: $({\rm S/N})^2 = \int_{r_{\min}}^{r_{\max}} n\,2\pi r\,\frac{A^2}{r^2\sigma^2}\,dr
= 2\pi nA^2\sigma^{-2}\ln(r_{\max}/r_{\min})$.
(b) $A^2/\sigma^2 = 0.0324$, so $2\pi\cdot5\cdot0.0324 = 1.02$, times $\ln 30 = 3.40$: $({\rm S/N})^2 = 3.46$, so
${\rm S/N} \approx 1.9$. (The actual galaxies give 1.65 for isotropic $\sigma = 0.049$.)
(c) Each factor of $e$ in radius contributes equally: within 1\,Mpc, $\ln 10/\ln 30 = 0.68$; within 0.5\,Mpc,
$\ln 5/\ln 30 = 0.47$. (The simulation gives 0.58 and 0.35; the NFW profile falls faster than $1/r$ at large $r$
and is shallower near the centre.) The signal is spread over a few Mpc, not concentrated in one spot, which is why a
method that uses every galaxy beats one cell.
(d) KL measures one shear direction much better than the other (median $\sigma \approx 0.039$ along the best direction
vs $0.057$ along the worst). With matrix weights, each galaxy contributes most where it is precise (A3). This is the
benchmark for Part G: no map can beat the matched filter.
}
